% Formal setting. Every number here comes from tables/numbers.tex, written by jr/tables.py
% from results/<cohort>/theory_numbers.json, and is checked by jr/theory_checks.py.

\section{Formal Setting}
\label{sec:theory}

This section fixes what ``no hindsight'' means for a Jira export and derives three consequences. Each one turns a pattern that is easy to miss into a prediction the experiments can check: how empty sprints inflate ROC-AUC, why a threshold fitted on one group of teams does not carry to another, and how the two spillover labels should differ.

\subsection{The event log and the sprint state}

A Jira export is an event log $\mathcal{E} = \{(i, f, v^{-}, v^{+}, t)\}$: issue $i$ had field $f$ changed from $v^{-}$ to $v^{+}$ at time $t$. With each issue's creation time $\gamma_i$ and resolution time $r_i$ (undefined while the issue is open), the log gives the value of any field at any time by replaying the changes in order. Write $\mathcal{E}_{\le t}$ for the events up to time $t$, and $\mathcal{F}_t$ for what they tell us together with the sprint calendar, which is fixed at planning. A sprint $s$ has a start $a_s$, a planned end $b_s$ and a closure time $c_s$.

The committed set is the board at the start,
\begin{equation}
C_s = \{\, i : s \in \mathrm{sprint}_i(a_s),\ \gamma_i \le a_s,\ r_i \not< a_s \,\},
\label{eq:committed}
\end{equation}
where $\mathrm{sprint}_i(t)$ is the issue's sprint field at time $t$ as replayed from $\mathcal{E}_{\le t}$. The outcomes are read off the end state. With $n_s = |C_s|$ and the carry ratio
\begin{equation}
\rho_s = \frac{|\{ i \in C_s : r_i > b_s \}|}{n_s},
\end{equation}
spillover is $y^{\mathrm{sp}}_s = \mathbf{1}[\rho_s > 0]$, serious spillover is $y^{\tau}_s = \mathbf{1}[\rho_s \ge \tau]$, late closure is $y^{\mathrm{late}}_s = \mathbf{1}[c_s - b_s > \delta]$ with $\delta$ one day (three for ``seriously late''), and scope churn is $\mathbf{1}[m_s / n_s \ge \theta]$, where $m_s$ counts issues added after $a_s$ and still in the sprint at $b_s$.

\subsection{Admissible features}

\begin{definition}[Admissible feature]
A feature $x(s, t)$ for sprint $s$ at prediction time $t$ is admissible if it is a function of $\mathcal{F}_t$ alone, that is, of $\mathcal{E}_{\le t}$ and the sprint calendar.
\end{definition}

Two common practices break this rule. First, a history feature that summarises an earlier sprint $u$ of the same project is admissible at $a_s$ only if $u$'s outcome was known by then, $\max(b_u, c_u) \le a_s$. A ``previous sprint'' picked by start date alone breaks this whenever sprints overlap, which happens for \OverlapShare\% of the sprints in our cohort. Second, a field read from the static export is $\mathcal{F}_\infty$-measurable, not $\mathcal{F}_{a_s}$-measurable. Wherever the log records no change for a field, the replay has to fall back on a value that may postdate planning. We record the share of committed issues affected, $\phi_s$, and repeat the main comparisons on sprints with $\phi_s = 0$. Issue type, priority and assignee are left out altogether, because the export only holds their final values.

\subsection{Empty sprints inflate ROC-AUC}

A sprint with $n_s = 0$ cannot spill, so it is a negative by construction. Keeping such sprints in a spillover evaluation rewards the model for something it did not have to learn.

\begin{proposition}
\label{prop:auc}
Split the negatives of a test set into structural negatives $N_0$ and ordinary negatives $N_1$, and let $\pi_0 = |N_0| / (|N_0| + |N_1|)$. Then
\begin{equation}
\mathrm{AUC}_{\mathrm{all}} = \pi_0\, A_0 + (1 - \pi_0)\, \mathrm{AUC}_{N_1},
\label{eq:aucmix}
\end{equation}
where $A_0$ is the probability that a random positive outscores a random structural negative, and $\mathrm{AUC}_{N_1}$ is the AUC on the positives and ordinary negatives alone. If the model ranks every structural negative below every positive, $A_0 = 1$ and the inflation is $\pi_0 (1 - \mathrm{AUC}_{N_1})$.
\end{proposition}

\begin{proof}
AUC is the probability that a random positive outscores a random negative. Conditioning on which kind of negative is drawn gives (\ref{eq:aucmix}).
\end{proof}

The decomposition is elementary. Its use is that it puts an exact number on the inflation.

On our time-split test set the full population has \NegAll{} negatives, and \NegStruct{} of them commit nothing, so $\pi_0 = \PiZero$. The model's AUC on nonempty sprints is \AucNonempty. With $A_0 = 1$, (\ref{eq:aucmix}) predicts \AucPred{} for the full population; the observed value is \AucAll, and the measured $A_0$ is \AZero. The gap of \Inflation{} between the two headline numbers is entirely the empty sprints. Any study that keeps them in the denominator reports this gap as skill.

\subsection{Why a fitted threshold does not travel}

Youden's index picks the threshold that maximises sensitivity plus specificity, which is the same as minimising $\mathrm{FNR} + \mathrm{FPR}$. With prevalence $\pi$, a miss cost $c_{\mathrm{FN}}$ and a false-alarm cost $c_{\mathrm{FP}}$, the expected cost per sprint is $\pi c_{\mathrm{FN}} \mathrm{FNR} + (1 - \pi) c_{\mathrm{FP}} \mathrm{FPR}$. This is proportional to $\mathrm{FNR} + \mathrm{FPR}$ exactly when
\begin{equation}
\frac{c_{\mathrm{FN}}}{c_{\mathrm{FP}}} = \frac{1 - \pi}{\pi}.
\label{eq:youden}
\end{equation}

\begin{proposition}
\label{prop:youden}
A Youden threshold fitted on a population with prevalence $\pi_{\mathrm{train}}$ is cost-optimal on a population with prevalence $\pi$ only if the implied cost ratio (\ref{eq:youden}) is the same, that is, only if $\pi = \pi_{\mathrm{train}}$. For a calibrated score the cost-optimal rule warns when $p \ge p^{*} = c_{\mathrm{FP}} / (c_{\mathrm{FP}} + c_{\mathrm{FN}})$, which does not depend on prevalence. A fixed score threshold carried to a team with lower prevalence therefore acts like a rule that prices a miss at only $(1 - \pi_{\mathrm{train}}) / \pi_{\mathrm{train}}$ false alarms, and its sensitivity falls with $\pi$.
\end{proposition}

The link between a threshold rule and the cost ratio it implies is standard in the ROC literature \cite{hernandez2012}. What we add is its use to explain why a threshold that worked on TAWOS misses most failures on some new teams.

For late closure the training prevalence is \TrainPrev, so the implied ratio is \TrainRatio. On the \ExtProjects{} testable outside projects the same threshold meets prevalences from \PrevMin{} to \PrevMax, where (\ref{eq:youden}) would call for ratios from \RatioMin{} down to \RatioMax. The observed sensitivity runs from \SensMin{} to \SensMax, with a rank correlation of \SpearmanPS{} with prevalence. The ranking travels because AUC does not depend on a threshold; the threshold does not travel. A deployment has to recalibrate it per team or set it from stated costs.

\subsection{A minimal model of the two spillover labels}

Suppose each committed issue of sprint $s$ is finished by the planned end independently with probability $q_s$, a property of the team and the period rather than of the issue. The number left unfinished is then $\mathrm{Binomial}(n_s, 1 - q_s)$, and
\begin{align}
\Pr[y^{\mathrm{sp}}_s = 1] &= 1 - q_s^{\,n_s}, \label{eq:any}\\
\mathbb{E}[\rho_s] &= 1 - q_s, \qquad \mathrm{Var}[\rho_s] = \frac{q_s (1 - q_s)}{n_s}. \label{eq:ratio}
\end{align}
Two predictions follow. By (\ref{eq:any}) the ``any'' label is driven by $n_s$ and saturates fast: at the cohort's median carry ratio, \CarryMedian, $q_s \approx \Qhat$ and a sprint of three issues already spills with probability \PSpillThree. The median nonempty sprint commits \MedianIssues{} issues, and \NonemptySpill\% of nonempty sprints spill. So committed scope should carry most of what can be known about plain spillover. By (\ref{eq:ratio}) a threshold label on $\rho_s$ is driven by $1 - q_s$, the team's completion rate, with $n_s$ entering only through the variance. Serious spillover should therefore depend on planning-time history and capacity rather than on scope, and the gap between scope-only and history-aware models should be larger for it than for plain spillover. Section~\ref{sec:results} tests both. The model ignores dependence between issues and scope changes during the sprint. It is a reading aid, not a claim about how teams work.

\subsection{Warnings as decisions}

A warning raised at checkpoint $c$ costs $\kappa$ and, if the sprint would otherwise fail, averts a fraction $e(c)$ of the failure cost $F$, with $e$ falling as $c$ grows. Warning at calibrated risk $p$ has expected cost $\kappa + F p (1 - e(c))$ against $F p$ for staying silent, so the warning pays when
\begin{equation}
p \ge \frac{\kappa}{e(c)\, F}.
\label{eq:warn}
\end{equation}
The bar rises as the sprint goes on, because a late warning saves less. An in-sprint model is useful only where its gain in ranking outpaces the fall in $e(c)$. Because (\ref{eq:warn}) needs a calibrated $p$, we report calibration beside AUC.
